\documentclass[11pt,a4paper]{article}

\usepackage{amsmath,amssymb}
\usepackage[utf8]{inputenc}
\usepackage{hyperref}
\usepackage{geometry}
\usepackage{booktabs}
\usepackage{xcolor}
\usepackage{enumitem}
\usepackage{setspace}

\geometry{margin=1in}
\onehalfspacing

\newcommand{\RH}{\ensuremath{\mathrm{RH}}}

\title{%
  Teaching Proof in the Age of Verification:\\
  Lessons from the Cathedral for Mathematics Education%
}
\author{
  Jason Robert Gochanour
}
\date{April 19, 2026}

\begin{document}

\maketitle

\begin{abstract}
The Cathedral---a machine-verified reduction of the Riemann Hypothesis
in Lean~4---was built in 23 days by a computer scientist (not a
professional mathematician) working with AI assistance. This paper
examines what this implies for mathematics education. We identify
five pedagogical shifts: (1)~proof as \emph{compilation} rather than
persuasion, (2)~error as \emph{data} rather than failure,
(3)~axioms as \emph{design choices} rather than sacred truths,
(4)~collaboration with AI as a \emph{cognitive amplifier} rather
than a crutch, and (5)~archival as \emph{curriculum} rather
than waste. We propose concrete changes to undergraduate and
graduate mathematics curricula, argue that formal verification
should be taught alongside traditional proof writing, and suggest
that the Cathedral's methodology---build, break, archive, audit---
is a teachable skill that transcends mathematics.
\end{abstract}

\tableofcontents

% ================================================================
\section{The Uncomfortable Question}
% ================================================================

A computer scientist with no PhD in mathematics, working with AI
assistance, produced a 78-file machine-verified proof system that
reduces one of the most famous open problems in mathematics to two
axioms. He did this in 23 days.

This raises questions that mathematics educators must confront:

\begin{enumerate}
  \item If the tools for formal verification are this accessible,
    should we be teaching proof differently?
  \item If AI can generate 85\% of proof code, what should humans
    learn?
  \item If machine verification can catch errors that experts miss
    (the High-Frequency Trap, the False Dedekind Reciprocity), what
    is the role of mathematical intuition?
  \item If 96 archived files (failed attempts) are essential to the
    discovery process, should we teach failure as curriculum?
\end{enumerate}

This paper does not claim that mathematicians are obsolete. The
Cathedral required deep mathematical taste, architectural vision,
and creative problem-solving that no AI currently possesses. But
the \emph{mechanics} of proof---the line-by-line construction of
logical arguments---have been fundamentally changed by verification
technology. Education must adapt.

% ================================================================
\section{Five Pedagogical Shifts}
% ================================================================

\subsection{Shift 1: Proof as Compilation}

\textbf{Traditional model:} A proof is a persuasive argument written
for a human reader. Its correctness is judged by whether experts
find it convincing.

\textbf{Cathedral model:} A proof is a program that compiles. Its
correctness is judged by whether the type checker accepts it. The
audience is the compiler, not a professor.

\textbf{Pedagogical implication:} Students can get \emph{immediate,
objective feedback} on whether their proof is correct. No waiting
for grading. No ambiguity about whether a step is justified. The
compiler either accepts it or it doesn't.

This changes the emotional experience of learning proof:

\begin{center}
\begin{tabular}{@{}lll@{}}
\toprule
& \textbf{Traditional} & \textbf{Verified} \\
\midrule
Feedback speed & Days (grading) & Seconds (compilation) \\
Feedback objectivity & Subjective & Objective \\
Error granularity & ``This step is wrong'' & Exact error location \\
Partial credit & Yes & No (but partial progress visible) \\
Student anxiety & High (uncertainty) & Lower (clarity) \\
\bottomrule
\end{tabular}
\end{center}

\textbf{Concrete proposal:} Introduce Lean~4 proof exercises alongside
traditional proof writing in discrete mathematics and introduction
to proof courses. Students write both a human-readable proof and a
machine-checked proof, learning that these are complementary skills.

\subsection{Shift 2: Error as Data}

\textbf{Traditional model:} Errors in proofs are failures. A wrong
proof is a bad grade.

\textbf{Cathedral model:} Errors are \emph{discoveries}. The
Cathedral's five traps---the High-Frequency Trap, the False Dedekind
Reciprocity, the Triangle Inequality Trap, the Hyperplane Trap, and
the Ghost Axiom Phenomenon---were each detected by the machine and
each redirected the proof toward a better path.

The 96 archived files are not failures. They are the empirical record
of a systematic exploration of mathematical space. The False Dedekind
Reciprocity, detected numerically at $(a,b) = (3,2)$ before any proof
attempt, saved potentially weeks of wasted effort and motivated the
Parseval Bridge---the key innovation.

\textbf{Pedagogical implication:} Students should be graded not just
on correct proofs but on their \emph{proof journals}: the record of
what they tried, what failed, and what they learned from failures.
The archive is as valuable as the proof.

\textbf{Concrete proposal:} In proof-based courses, require students
to submit a ``proof diary'' alongside each assignment: what approaches
they tried, why each failed, and how failures informed their
final approach. Grade the diary.

\subsection{Shift 3: Axioms as Design Choices}

\textbf{Traditional model:} Axioms are foundational truths that
mathematicians agree on. Students learn to prove things \emph{from}
axioms but rarely think about the axioms themselves.

\textbf{Cathedral model:} Axioms are \emph{engineering decisions}.
The Cathedral started with 56 axioms and reduced to 39 (7 on the
crown path) through systematic campaigns. Each axiom was named,
classified by tier, tracked by dependents, and evaluated for
elimination. This is axiom \emph{management}, not axiom worship.

\textbf{Pedagogical implication:} Students should learn to:
\begin{enumerate}
  \item Identify the axioms in a proof (what are we assuming?).
  \item Classify axioms by strength (how much does this assumption
    buy us?).
  \item Ask whether axioms can be eliminated (can we prove this
    from something weaker?).
  \item Understand the trade-off between axiom count and proof
    complexity.
\end{enumerate}

\textbf{Concrete proposal:} Add ``axiom awareness'' to proof
courses. For each proof, students must list the axioms used and
classify them as ``proved in Mathlib,'' ``standard but unformalized,''
or ``genuinely new assumption.''

\subsection{Shift 4: AI as Cognitive Amplifier}

\textbf{Traditional model:} Students must produce proofs
independently. Using tools beyond pen and paper is cheating.

\textbf{Cathedral model:} AI generated 85\% of code by volume.
The human provided 100\% of the architectural vision. The result
was a system that neither could have produced alone.

\textbf{Pedagogical implication:} We need to teach students how
to \emph{collaborate with AI} on mathematical reasoning:

\begin{itemize}
  \item How to decompose a problem into tasks suitable for AI
    assistance.
  \item How to evaluate AI-generated proofs (the compiler helps,
    but understanding the proof is still the human's job).
  \item How to provide architectural direction that AI cannot
    generate on its own.
  \item When to trust AI output and when to verify independently.
\end{itemize}

The analogy: we teach students to use calculators without
considering it cheating. We should teach them to use proof
assistants and AI similarly---as tools that handle mechanical
computation so the human can focus on creative reasoning.

\textbf{Concrete proposal:} Develop ``AI-assisted proof'' exercises
where students are given a theorem and must:
\begin{enumerate}
  \item Formulate a proof strategy (human).
  \item Use AI to generate candidate proof steps (AI).
  \item Verify the result in Lean~4 (compiler).
  \item Write a human-readable explanation (human).
\end{enumerate}

\subsection{Shift 5: Archive as Curriculum}

\textbf{Traditional model:} Published proofs are clean, polished
artifacts. The messy process of discovery is hidden.

\textbf{Cathedral model:} The archive (96 files) is preserved
alongside the active codebase (78 files). Failed approaches are
documented with explanations of \emph{why} they failed. The
archive is a teaching resource.

\textbf{Pedagogical implication:} Students should study
\emph{failed proofs} as seriously as successful ones. Understanding
why an approach fails develops mathematical maturity faster than
reading correct proofs.

\textbf{Concrete proposal:} In graduate courses on analytic number
theory or proof engineering, assign the Cathedral's archive as
reading material. Students analyze a failed approach (e.g., the
Vasyunin cotangent path) and write a report on:
\begin{enumerate}
  \item What the approach attempted.
  \item Where and why it failed.
  \item What was learned from the failure.
  \item How the failure informed the eventual solution.
\end{enumerate}

% ================================================================
\section{Curriculum Recommendations}
% ================================================================

\subsection{Undergraduate}

\begin{enumerate}
  \item \textbf{Introduction to Proof (Year 1--2):} Add Lean~4
    exercises alongside traditional proofs. Start with propositional
    logic and basic set theory in Lean.

  \item \textbf{Discrete Mathematics (Year 1--2):} Use Lean~4 for
    induction proofs, divisibility, and modular arithmetic.
    Students see that the computer checks their induction step
    the same way a professor would.

  \item \textbf{Linear Algebra (Year 2):} Formalize basic matrix
    theorems (determinant properties, eigenvalue bounds). The
    Cathedral's Sherman--Morrison and Schur Complement proofs
    are accessible at this level.

  \item \textbf{Analysis (Year 2--3):} Use Lean~4 to formalize
    $\varepsilon$-$\delta$ proofs. The precision required by the
    compiler teaches students to be exact about quantifiers.
\end{enumerate}

\subsection{Graduate}

\begin{enumerate}
  \item \textbf{Formal Methods seminar:} Study the Cathedral as
    a case study in large-scale formalization. Topics:
    axiom management, proof architecture, the archive as
    documentation.

  \item \textbf{Analytic Number Theory:} Use the Cathedral's
    proof tree as a roadmap. Students attempt to formalize
    one of the 38 non-critical axioms as a class project.

  \item \textbf{AI-Assisted Mathematics:} Study the human-AI
    collaboration model. Compare with AlphaProof, LeanDojo,
    and other AI-for-math systems. Discuss the ``cognitive
    amplifier'' vs ``replacement'' debate.
\end{enumerate}

% ================================================================
\section{What Students Should Learn That Won't Change}
% ================================================================

Machine verification does not replace mathematical understanding.
The Cathedral required:

\begin{itemize}
  \item \textbf{Taste}: Recognizing that the Parseval Bridge was
    more promising than the Vasyunin path. No AI suggested this.
  \item \textbf{Intuition}: Sensing that the Rank-1 Mellin Miracle
    ``should'' work before proving it formally.
  \item \textbf{Architecture}: Designing the 11-module structure
    that made the proof manageable. This is software engineering
    applied to mathematics.
  \item \textbf{Persistence}: 96 failed files. Three weeks of
    12-hour days. The machine checks the proof; the human must
    keep going when everything is failing.
  \item \textbf{Communication}: 17 papers for 17 audiences. The
    ability to explain the same idea to a physicist, a philosopher,
    and a senator. No compiler checks this.
\end{itemize}

These are human skills. They should remain at the center of
mathematics education. What changes is the \emph{mechanics}:
the laborious, error-prone process of checking logical steps by
hand is increasingly delegated to machines, freeing human minds
for the creative work that machines cannot yet do.

% ================================================================
\section{Conclusion: The Compiler as Teacher}
% ================================================================

The Lean compiler is, in a sense, the most patient and honest
mathematics teacher imaginable. It never gets tired. It never
gives partial credit for a wrong proof. It never misses an error.
It gives feedback in seconds. And it treats every student exactly
the same way, regardless of background, reputation, or pedigree.

A computer scientist in New Mexico, with no doctorate in
mathematics, working with AI tools available to anyone, built a
78-file formal proof system that addresses the most famous open
problem in mathematics. He did this because the tools are now
good enough that mathematical \emph{taste} and \emph{persistence}
matter more than mathematical \emph{pedigree}.

If we teach the next generation to use these tools---not as
replacements for thinking, but as amplifiers for it---the next
cathedral will be built by a student who hasn't been born yet.

\begin{quote}
\emph{The compiler does not care who you are.\\
It only cares whether you are right.\\
That is the most democratic thing\\
in all of mathematics.}
\end{quote}

\begin{thebibliography}{99}

\bibitem{lean4}
L.~de~Moura \emph{et al.}, \emph{The Lean~4 Theorem Prover and
Programming Language}, CADE-28, 2021.

\bibitem{mathlib}
The Mathlib Community, \emph{Mathlib: A unified library of
mathematics formalized in Lean~4}, 2024.

\bibitem{avigad}
J.~Avigad, \emph{Mathematics in Lean}, online textbook, 2023.

\bibitem{buzzard}
K.~Buzzard, \emph{The Future of Mathematics}, ICM 2022 invited
lecture.

\bibitem{hales}
T.~Hales, \emph{Formal Abstracts}, Fields Institute, 2019.

\end{thebibliography}

\end{document}
